\documentclass[conference]{IEEEtran}
\IEEEoverridecommandlockouts
\usepackage{cite}
\usepackage{amsmath,amssymb,amsfonts}
\usepackage{graphicx}
\usepackage{textcomp}
\usepackage{xcolor}
\usepackage{booktabs}
\usepackage{tabularx}
\usepackage{array}
\usepackage{balance}
\def\BibTeX{{\rm B\kern-.05em{\sc i\kern-.025em b}\kern-.08em\TeX}}
\begin{document}

\title{Can We Cure Neurodegenerative Diseases?\\
Stage-Adaptive Proteostasis--Immune--Repair Research Platform}

\author{\IEEEauthorblockN{Source-Controlled Research Plan}
\IEEEauthorblockA{Survey--Idea--ExperimentDesign--Author pipeline\\
Translational neuroscience and biomedical mechanisms}}

\maketitle

\begin{abstract}
Neurodegenerative diseases cause progressive neuronal and circuit failure, yet the word ``cure'' is often used for outcomes that differ fundamentally: symptomatic relief, slowing of decline, pathological remission, and repair of lost function. This paper proposes a Stage-Adaptive Proteostasis--Immune--Repair Platform (SPIRP) for Alzheimer disease (AD) and Parkinson disease (PD). SPIRP combines disease-specific protein modules with a shared systems-level backbone. Patient-derived induced pluripotent stem cell (iPSC) neurons and glia are used for responder selection and target engagement; staged animal models test causal treatment order; and longitudinal biomarker cohorts test whether a composite state classifier generalizes to human disease. The design measures aggregate burden, clearance capacity, glial state, neuronal reserve, circuit function, toxicity, and rebound after treatment withdrawal. It explicitly compares proteostasis restoration, immune-state normalization, mitochondrial or synaptic resilience, and conditional stem-cell-derived neural repair. The main hypothesis is that durable benefit requires a state-matched sequence: pathology reduction and clearance restoration, followed by selective immune normalization and only then repair when the tissue niche is permissive. A treatment is not called curative unless functional control, biological remission, safety, durability, and independent replication are all demonstrated. The report is a research proposal with no newly collected data and no clinical treatment recommendation.
\end{abstract}

\begin{IEEEkeywords}
neurodegeneration, Alzheimer disease, Parkinson disease, proteostasis, neuroinflammation, microglia, drug repurposing, induced pluripotent stem cells, neural repair, disease modification
\end{IEEEkeywords}

\section{Introduction}
Neurodegenerative diseases are not one disease with one molecular cause. Alzheimer disease and Parkinson disease differ in their characteristic protein pathologies, vulnerable neuronal populations, circuit failures, clinical trajectories, and genetic and environmental modifiers. They nevertheless share several biological processes: misfolding and aggregation of proteins, impaired degradation and trafficking, synaptic and axonal dysfunction, mitochondrial stress, glial responses, and progressive loss of neuronal reserve \cite{selkoe2016,poewe2017,bloem2021}. This combination raises an attractive but dangerous possibility. A common therapeutic architecture may exist, but a common drug or universal cure cannot be assumed from convergent mechanisms alone.

The question ``Can we cure neurodegenerative diseases?'' is therefore reformulated as a sequence and state problem. Can a biomarker-stratified intervention first reduce toxic protein burden and restore clearance, then normalize maladaptive immune states while preserving protective functions, and finally protect or replace vulnerable neurons? Can the resulting functional benefit persist after treatment withdrawal? Which parts of that sequence apply to AD and PD, and which must remain disease-specific?

This paper presents the Stage-Adaptive Proteostasis--Immune--Repair Platform (SPIRP). It is the output of a four-stage evidence-to-design process. Survey identifies a gap ledger; Idea selects a mechanistic direction; ExperimentDesign specifies causal tests; and Author binds the proposal to a fixed evidence registry. The project is limited to molecular neuroscience, neuroimmunology, translational pharmacology, disease modeling, and neural repair. It does not analyze neurodegeneration through non-biological frameworks, and it does not report a completed experiment.

SPIRP makes a stronger claim than a conservative inventory of limitations. It proposes that some forms of neurodegeneration may become durably controllable, and that repair may be feasible for a subset of patients if intervention is matched to molecular state and tissue reserve. It also makes that ambition falsifiable. A treatment that improves a symptom without changing pathology or durability is classified as symptom control. A treatment that clears pathology but does not restore function is classified as mechanism-only. A graft that survives but fails to integrate is not neural repair. The platform is designed to discover these distinctions rather than hide them inside a single endpoint.

\section{Background, Survey, and Research Gap}
\subsection{Convergent mechanisms do not imply one disease}
The amyloid hypothesis of AD provides a mechanistic account in which amyloid-beta accumulation can initiate or amplify tau pathology, synaptic injury, inflammation, and neuronal dysfunction \cite{selkoe2016}. It does not imply that every patient has the same initiating event, nor that removing amyloid reverses downstream neuronal loss. PD involves alpha-synuclein aggregation and Lewy pathology, dopaminergic vulnerability in the substantia nigra, and a broad set of motor and non-motor circuit processes \cite{poewe2017,bloem2021}. A platform should therefore share methods for measuring proteostasis and reserve, while using amyloid/tau and alpha-synuclein/nigrostriatal modules as distinct disease components.

The practical implication is a modular causal graph. Protein burden affects synapses and organelles; clearance capacity determines whether burden resolves; glial state modulates debris handling and inflammatory injury; neuronal reserve determines recoverability; and circuit function integrates the resulting tissue state. Age, vascular state, metabolism, genotype, sex, medication, and systemic inflammation modify several edges at once. This graph is more useful for treatment selection than the label ``neurodegenerative disease'' alone.

\subsection{Protein clearance and clinical disease modification}
Protein aggregation is a target, but it is also a timing problem. In early AD, the lecanemab trial demonstrated a slowing of clinical decline alongside amyloid reduction in a selected population \cite{vandyk2023}. This is important evidence that pathology-directed treatment can modify disease trajectory. It is not evidence of complete cure, neuronal replacement, or universal efficacy. Treatment after extensive neuronal loss may remove amyloid while leaving irreversible circuit damage. SPIRP therefore requires a target-engagement measure, a neuronal-reserve estimate, and a washout period.

For PD, alpha-synuclein biology is not interchangeable with amyloid biology. A common proteostasis backbone can test lysosomal and autophagic capacity, mitochondrial stress, or aggregate clearance, but the disease module must retain alpha-synuclein burden, dopaminergic neuron survival, and striatal circuit function. The platform deliberately preserves this difference instead of using a generic aggregate score as a substitute for pathology.

\subsection{Neuroinflammation and glial state}
Neuroinflammation is a central feature of AD biology, but it is not a single harmful switch. Microglia and astrocytes can clear debris, remodel synapses, respond to injury, and support repair; persistent or maladaptive states can amplify synapse loss and neuronal stress. Heneka et al. described neuroinflammation as a major component of AD pathogenesis while emphasizing its complex biology \cite{heneka2015}. Single-cell profiling identified disease-associated microglial states rather than one uniform activated population \cite{keren2017}. Broad immune suppression could therefore remove protective activity together with harmful signaling.

SPIRP uses state-selective immune normalization as an experimental objective. It measures clearance, inflammatory outputs, phagolysosomal function, synaptic interaction, and tissue toxicity separately. A reduction in a cytokine panel is not sufficient. The desired state is one that reduces maladaptive neuronal injury while preserving or improving aggregate handling and repair compatibility.

\subsection{Human cellular models and their boundaries}
Human iPSC models can preserve patient genetic backgrounds and permit direct study of neurons, astrocytes, and microglia. Reviews of iPSC models in neurodegeneration show their utility for studying autophagy and disease-linked cellular phenotypes \cite{seranova2020}. They are particularly valuable for comparing responders and non-responders before animal studies. However, they do not automatically reproduce biological aging, vascular perfusion, systemic immunity, long-range brain connectivity, or the full time course of disease. SPIRP uses iPSC models as a human molecular testbed, not as a complete substitute for patients.

\subsection{Drug repurposing and neural repair}
Drug repurposing can accelerate candidate selection because pharmacokinetic and pharmacodynamic information may already exist. Yet prior approval for one indication does not establish brain exposure, target engagement, combination compatibility, or efficacy in a neurodegenerative stage. Pushpakom et al. describe repurposing as a strategy with substantial promise but also biological, regulatory, and evidence-quality obstacles \cite{pushpakom2019}. SPIRP treats repurposing as a mechanism-first discovery funnel.

Stem-cell-derived neuronal or glial replacement addresses a different problem: lost cells and circuit function. It cannot be expected to solve an unresolved pathological environment. Graft survival, phenotypic maturity, synaptic integration, immune compatibility, tumorigenicity, ectopic activity, and long-term pathology are required endpoints. Neural repair is consequently a late or conditional module, not a standalone cure claim.

\subsection{Gap ledger}
The Survey stage accepted eight gaps. G1 is the absence of a unified cure endpoint. G2 is the missing link between molecular stage and treatment response. G3 concerns the order and interaction of proteostasis, immune, resilience, and repair interventions. G4 is the lack of a unified biomarker framework for selection and target engagement. G5 is the incomplete representation of aging and whole-brain disease in iPSC and organoid systems. G6 is the missing validation of CNS exposure and stage-specific efficacy for repurposed drugs. G7 is the unresolved requirement for a permissive niche before graft repair. G8 is the limited testing of durability after treatment withdrawal.

\begin{table*}[t]
\caption{Evidence anchors and inference boundaries}
\label{tab:evidence}
\centering
\small
\begin{tabularx}{\textwidth}{p{0.08\textwidth}p{0.32\textwidth}p{0.27\textwidth}X}
\toprule
ID & Anchor & Contribution & Boundary \\
\midrule
E1 & Selkoe and Hardy, 2016 \cite{selkoe2016} & Amyloid, downstream injury, and AD mechanism & Does not establish a universal cure or reversal \\
E2 & Poewe et al., 2017 \cite{poewe2017} & PD pathology, progression, and therapeutic context & PD remains biologically heterogeneous \\
E3 & Bloem et al., 2021 \cite{bloem2021} & PD clinical and circuit overview & Clinical heterogeneity limits one endpoint \\
E4 & Heneka et al., 2015 \cite{heneka2015} & AD neuroinflammation and immune mechanisms & Immune effects are state- and stage-dependent \\
E5 & Keren-Shaul et al., 2017 \cite{keren2017} & Disease-associated microglial states & No single microglial state represents all disease \\
E6 & van Dyck et al., 2023 \cite{vandyk2023} & Early AD anti-amyloid disease modification & Slowing is not cure or neuronal repair \\
E7 & Seranova et al., 2020 \cite{seranova2020} & Human iPSC models and autophagy & Cellular models incompletely represent aging and circuits \\
E8 & Pushpakom et al., 2019 \cite{pushpakom2019} & Drug repurposing opportunities and obstacles & Existing approval does not prove CNS efficacy \\
\bottomrule
\end{tabularx}
\end{table*}

\section{Research Questions and Planned Contributions}
SPIRP asks four questions. First, can composite biomarkers identify a state in which proteostasis restoration changes pathology and neuronal reserve? Second, does selective immune-state normalization improve neuronal survival more effectively than broad suppression? Third, does the order of pathology control, immune normalization, and neural repair determine durability? Fourth, can a common systems backbone generalize across AD and PD while preserving disease-specific protein and circuit modules?

The planned contributions are: (1) a nested operational definition of functional control, biological remission, and repair; (2) a state-space model that joins aggregate burden, clearance capacity, glial state, neuronal reserve, and circuit function; (3) a factorial sequence experiment spanning human cellular models and animal disease modules; (4) a washout-centered durability endpoint; and (5) an explicit rule for accepting or rejecting cross-disease generalization.

The proposal is deliberately ambitious. It does not assume that current therapies are cures, but it does not stop at symptom management. It creates a route by which a true durable remission, if it exists for a defined molecular state, can be distinguished from short-term improvement. Its strongest possible outcome is not a single universal treatment. It is a validated state-to-sequence map that tells researchers when pathology clearance, immune normalization, or repair is sufficient and when combinations are required.

\section{Idea Source Checkpoints and Direction Selection Audit}
The Idea stage compared five routes. A universal anti-aggregation drug was rejected because AD and PD have disease-specific proteins and downstream circuits. Broad anti-inflammatory treatment was rejected as a primary route because protective and maladaptive glial functions are not equivalent. A repurposing-only screen was retained as a candidate discovery engine but lacks durability and repair. A repair-first strategy was rejected because unresolved pathology can damage transplanted cells and prevent integration. SPIRP was selected because it covers all eight gaps and can be staged from human cellular tests to causal animal studies and biomarker validation.

The evolution had four checkpoints. First, cure was changed from a binary label to nested functional, remission, and repair endpoints. Second, the single mechanism was changed to a shared systems backbone with disease-specific pathology modules. Third, cell replacement was moved from the beginning to a conditional stage after pathology and immune-state control. Fourth, a single biomarker was replaced by a locked composite state score and held-out prediction. These changes increase falsifiability and make the proposed sequence more informative than an unstructured combination trial.

\begin{table}[t]
\caption{Idea route audit}
\label{tab:audit}
\centering
\scriptsize
\setlength{\tabcolsep}{2pt}
\resizebox{\columnwidth}{!}{%
\begin{tabular}{lccccc}
\toprule
Route & Novelty & Causal & Translation & Gaps & Decision \\
\midrule
Universal anti-aggregate & Low & Medium & Medium & 3/8 & Reject \\
Broad immune suppression & Medium & Medium & Medium & 3/8 & Reject \\
Repurposing screen only & High & Low & High & 5/8 & Component \\
Repair first & High & Medium & Medium & 4/8 & Reject \\
SPIRP & High & High & High & 8/8 & Primary \\
\bottomrule
\end{tabular}%
}
\end{table}

\section{Problem Definition, Assumptions, and Hypotheses}
The unit of analysis is a molecularly characterized disease state, not a diagnostic label alone. The state vector for individual or culture (i) at time (t) is
\begin{equation}
\mathbf{x}_{it}=[A_{it},T_{it},C_{it},M_{it},N_{it},V_{it}]^{\mathsf T},
\label{eq:state}
\end{equation}
where (A) is toxic aggregate burden, (T) is tau or alpha-synuclein propagation burden, (C) is clearance and proteostasis capacity, (M) is glial inflammatory state, (N) is neuronal reserve, and (V) is circuit function. The disease module determines which protein burden is emphasized. Age-equivalent maturation, genotype, sex, vascular/metabolic state, and batch are covariates rather than hidden assumptions.

Under treatment input \(\mathbf{u}_{it}\), the state follows the phase-indexed model
\begin{equation}
\mathbf{x}_{i,t+1}=\mathbf{F}_{\phi}(\mathbf{x}_{it},\mathbf{u}_{it},\mathbf{h}_{it})+\boldsymbol{\epsilon}_{it},
\label{eq:dynamics}
\end{equation}
where \(\phi\) is molecular stage, \(\mathbf{h}\) contains pharmacodynamic and systemic variables, and \(\boldsymbol{\epsilon}\) represents process variation. Measurements are represented by
\begin{equation}
\mathbf{y}_{it}=\mathbf{H}_{\phi}\mathbf{x}_{it}+\mathbf{D}\mathbf{z}_{it}+\boldsymbol{\eta}_{it},
\label{eq:observation}
\end{equation}
where \(\mathbf{z}\) includes assay batch, motion, locomotion, and other nuisance variables.

The composite state score is
\begin{equation}
S_{it}=w_A\tilde{A}_{it}+w_T\tilde{T}_{it}-w_C\tilde{C}_{it}+w_M\tilde{M}_{it}-w_N\tilde{N}_{it}-w_V\tilde{V}_{it},
\label{eq:score}
\end{equation}
where direction-aligned standardized features are marked by tildes and weights are learned in training data only. A higher score indicates greater vulnerability. The classifier is useful only if it predicts held-out response and remains calibrated across donors, animals, batches, and disease modules.

The hypotheses are: \textbf{H1}, proteostasis intervention improves aggregate clearance and neuronal resilience only with validated target engagement; \textbf{H2}, state-selective immune normalization improves neuronal survival while preserving protective clearance; \textbf{H3}, pathology-first or pathology-plus-immune-first sequences outperform repair-first sequences for durable graft survival and circuit rescue; \textbf{H4}, the composite score predicts held-out response better than any single biomarker; \textbf{H5}, remission requires low rebound after washout; and \textbf{H6}, a shared systems backbone generalizes directionally across AD and PD while protein and circuit modules remain disease-specific.

\section{Study Design and Methods}
\subsection{Human patient-derived cellular stage}
The cellular panel uses at least three independent donor backgrounds per disease module, with isogenic correction or introduction of selected variants where technically justified. It includes excitatory and dopaminergic neurons, astrocytes, and microglia. AD cultures emphasize amyloid/tau-linked phenotypes; PD cultures emphasize alpha-synuclein and dopaminergic vulnerability. Cell identity, maturation, neuronal activity, glial composition, genomic stability, and differentiation batch are quality-control variables.

The intervention library is organized by mechanism: autophagy-lysosomal or other proteostasis enhancement; disease-protein production, aggregation, or clearance modulation; microglial-state and inflammasome-pathway modulation; mitochondrial and oxidative-stress support; trophic and synaptic resilience; and stem-cell-derived neuronal or glial repair. Repurposed compounds are prioritized when CNS exposure and pharmacodynamic measurements are plausible. Novel compounds are retained when they offer a distinct target-engagement mechanism. Each candidate is screened for cytotoxicity, exposure range, interaction, and off-target effects before combinations.

For each module, cultures are randomized to vehicle, proteostasis alone, immune-state alone, simultaneous combination, proteostasis-first, immune-first, and the leading sequences followed by a repair module. Repair is added only after a predefined reduction in aggregate burden and maladaptive inflammatory-state score. No-pathology controls identify direct toxicity. Primary endpoints are aggregate load, validated proteostasis flux, neuronal survival, synaptic density or function, microglial state, and neuronal activity. Secondary endpoints include transcriptomic state, mitochondrial function, secreted inflammatory mediators, and cell-cell interaction.

The sequence estimand is
\begin{equation}
\Delta_{q,\phi}=E[Y\mid do(Q=q),\phi]-E[Y\mid do(Q=vehicle),\phi],
\label{eq:sequence}
\end{equation}
where $Q$ denotes an intervention sequence and $\phi$ the molecular stage. The sequence is promoted only if its intended endpoint improves without loss of neuronal viability or excessive suppression of protective glial functions.

\subsection{In vivo causal stage}
The AD module uses a validated amyloid/tau model with pathology, synaptic, circuit, and cognitive endpoints. The PD module uses a validated alpha-synuclein or dopaminergic-neurodegeneration model with alpha-synuclein, nigrostriatal, motor, and neuronal-survival endpoints. Models are locked before treatment based on phenotype stability and external replication. Animals are randomized within sex, litter, cage, age, and baseline severity; treatment allocation, behavior scoring, tissue processing, and image analysis are blinded.

Four stage strata are tested: pre-pathology risk, early measurable pathology with reserve, established pathology with partial reserve, and advanced loss. The objective is to estimate where a sequence can still achieve remission and where repair is required. Treatment arms compare single modules, simultaneous combination, pathology-first, immune-first, pathology-plus-immune followed by repair, and repair-first as a negative ordering control. Vehicle, sham, and non-disease controls are included. Treatment timing and washout are prespecified.

The repair module uses characterized stem-cell-derived neuronal or glial populations only after identity, purity, genomic stability, tumorigenicity, and dose-range requirements pass review. It measures graft survival, phenotypic maturity, synaptic integration, circuit activity, ectopic activity, immune response, tumor or overgrowth signals, and recurrent pathology. The design does not presume that a surviving graft is a functional graft.

Across AD and PD, primary endpoints include disease-specific pathology, neuronal survival, and circuit function. AD measures include amyloid/tau burden, synaptic density, hippocampal or cortical circuit function, and cognitive task performance. PD measures include alpha-synuclein burden, substantia nigra dopaminergic neuron counts, striatal dopamine function, and motor performance. Shared endpoints include glial state, lysosomal/proteostasis markers, systemic inflammation, body weight, general health, adverse events, and post-treatment rebound.

The rebound endpoint is
\begin{equation}
R_m=Y_{m,\mathrm{washout}}-Y_{m,\mathrm{end}},
\label{eq:rebound}
\end{equation}
with direction conventions fixed for each measure. Favorable remission requires sustained functional benefit and no predefined pathological rebound. A transient end-of-treatment improvement is classified as control, not remission.

\subsection{Translational biomarker stage}
The translational stage uses longitudinal observational cohorts spanning biomarker-defined risk, early disease, established disease, and matched controls. Blood and, where clinically justified, CSF, molecular imaging, structural or functional imaging, digital motor or cognitive measures, medication, vascular, and systemic covariates are collected. The goal is to validate whether $S_{it}$ predicts progression and target engagement in held-out participants. This stage is not a clinical efficacy trial.

If a later interventional trial is proposed, it requires a new protocol with regulatory review, eligibility criteria, dose and interaction analysis, stopping rules, adverse-event monitoring, and a prespecified clinical endpoint package. The present design does not authorize off-label treatment, experimental transplantation, or personal biomarker interpretation.

\subsection{Statistical analysis}
Cellular data use donor-aware mixed-effects models with donor, differentiation batch, and plate as random effects. Animal data use longitudinal mixed-effects models with sequence, disease stage, sex, baseline severity, and time; cage and litter are random effects. Translational data use joint models of biomarker trajectory and progression. The primary estimands are sequence-by-stage interactions and washout rebound rather than isolated within-group significance.

Power is determined separately for donor replication, animal stage-by-sequence interaction, and translational prediction using pilot variance, minimum meaningful difference, attrition, and multiplicity. Grouped cross-validation keeps all observations from one donor, animal, participant, or batch in one fold. The state classifier must improve held-out prediction over a simple disease-stage baseline and remain calibrated across disease modules. Missingness, assay limits, motion thresholds, exclusions, and model versions are fixed before unblinding.

\section{Expected Outcome Branches and Conditional Conclusions}
In Branch A, proteostasis treatment improves clearance and neuronal reserve with verified target engagement. This supports a state-specific mechanism and justifies immune or repair sequencing. In Branch B, selective immune normalization reduces neuronal injury while preserving phagolysosomal and clearance functions. This supports a glial-state mechanism, not broad immune suppression. In Branch C, pathology-first or pathology-plus-immune-first treatment improves graft integration and post-washout circuit function relative to repair-first treatment. This supports a permissive-niche requirement for neural repair.

In Branch D, function improves without pathology or durability improvement. The intervention is classified as symptom-control and not disease modification. In Branch E, pathology changes but neuronal function does not. The mechanism is retained for combination or earlier-stage testing, but no cure claim is made. In Branch F, the composite biomarker fails to predict held-out response. The state classifier is rejected and the biomarker strategy must be revised. In Branch G, the AD and PD sequences do not generalize. The shared backbone claim is rejected while disease-specific programs remain valid.

The conditional conclusion is direct. If target engagement, pathology reduction, functional recovery, safety, and low washout rebound replicate independently, a defined disease stage may meet the operational standard for biological remission. If neuronal repair adds durable circuit recovery in a controlled niche, a subset of disease states may meet the stronger repair standard. If these criteria are not met, the correct result is not that cure is impossible; it is that the tested sequence, state definition, or disease module is insufficient. This distinction keeps the research ambitious without converting preliminary promise into a clinical guarantee.

\section{Risks, Limitations, and Human Review Requirements}
The first risk is endpoint inflation. A slower decline is not a cure, and an imaging change is not necessarily functional recovery. The nested endpoint hierarchy and washout are designed to prevent these substitutions. The second risk is model incompleteness. iPSC cultures lack complete aging and circuits; animal models do not reproduce all human heterogeneity. Agreement across models is therefore treated as convergence of mechanisms, not proof of clinical efficacy.

The third risk is immune misinterpretation. A lower inflammatory marker may reflect loss of protective microglia rather than improvement. State-resolved transcriptional, phagolysosomal, neuronal, and toxicity measures are required. The fourth risk is treatment-order confounding. Sequence arms must be dose-, exposure-, and time-matched so that a later arm is not disadvantaged merely by receiving less total treatment or a different observation window.

The fifth risk is repurposing overconfidence. Established pharmacokinetics do not guarantee brain exposure in a new population or compatibility with a combination. CNS exposure, target engagement, interaction, and adverse-event monitoring remain mandatory. The sixth risk is repair toxicity. Stem-cell-derived grafts can fail to mature, integrate ectopically, trigger immune responses, or develop overgrowth. Identity, genomic stability, tumorigenicity, circuit activity, and long-term surveillance are release criteria.

The seventh risk is biomarker leakage. If samples from the same donor, animal, or participant occur in both training and test folds, apparent prediction can be inflated. Grouped cross-validation and external replication are mandatory. The eighth risk is population stratification. Genotype and biomarker models require ancestry-aware analysis, sex and age reporting, medication and vascular covariates, and privacy-preserving data governance. No molecular score may be used to label an individual as curable or incurable.

Animal work requires institutional animal-care review, 3R justification, minimized distress, humane endpoints, and transparent exclusion reporting. Human biomarker studies require informed consent, secure health-data handling, clinically justified sampling, and clear communication that the classifier is investigational. Any later clinical combination or transplantation study requires a new regulatory and ethics review. The platform does not recommend self-medication or unapproved cell therapy.

\section{Conclusion}
Can we cure neurodegenerative diseases? The evidence supports neither a universal no nor an immediate yes. It supports a more useful proposition: durable remission and, for selected states, neural repair may be reachable only through stage-adaptive treatment that couples disease-protein clearance, proteostasis restoration, state-selective immune normalization, and controlled circuit repair. SPIRP turns that proposition into a causal research program.

Its decisive tests are not a single biomarker or a short-term symptom score. They are verified target engagement, disease-specific pathology reduction, preserved or restored neuronal reserve, circuit function, low rebound after washout, safety, and independent replication. A positive result would identify a biological state and sequence in which cure-like outcomes are attainable. A negative result would identify which shared-backbone or repair assumptions fail. In either case, the platform replaces an imprecise promise with a rigorous map of what can be controlled, what can be repaired, and what remains unresolved.

\appendices
\section{Variables and Operational Definitions Appendix}
\begin{table*}[h]
\caption{Operational variables inherited across Survey, Idea, ExperimentDesign, and Author}
\label{tab:variables}
\centering
\small
\begin{tabularx}{\textwidth}{p{0.19\textwidth}p{0.32\textwidth}X}
\toprule
Variable & Operational definition & Interpretation rule \\
\midrule
Functional control & Slowed or arrested progression on prespecified function and circuit endpoints during treatment & Not equivalent to cure \\
Biological remission & Pathology and maladaptive immune state remain improved through washout without rebound & Requires target engagement and replication \\
Repair & Durable circuit recovery with safe cell survival and integration after pathology control & Requires long-term safety and function \\
Aggregate burden & Disease-specific amyloid/tau or alpha-synuclein burden by validated assay & Protein modules remain disease-specific \\
Proteostasis capacity & Flux or clearance assay for autophagy-lysosomal and related degradation function & Must be measured, not inferred from a marker \\
Glial state & Multi-feature state including inflammatory outputs, phagolysosomal function, and neuronal interaction & No uniform activated-state assumption \\
Neuronal reserve & Surviving neurons, synapses, activity, and recoverability measures & Determines repair opportunity \\
Circuit function & Cognitive, motor, electrophysiological, imaging, or synaptic function appropriate to model & Must be interpreted with pathology \\
Composite score & Locked state score trained on orthogonal features & Requires held-out calibration \\
Washout rebound & Difference between washout and end-of-treatment measures & High rebound defeats remission claim \\
\bottomrule
\end{tabularx}
\end{table*}

\section{Evidence Coverage and Review Appendix}
The Survey evidence plan assigns E1--E3 to convergent and disease-specific mechanisms; E4--E5 to neuroinflammation and microglial state; E6 to the boundary between disease modification and cure; E7 to human cellular modeling and autophagy; and E8 to drug repurposing. All eight accepted gaps are inherited by SPIRP. The primary direction is linked to each gap. The document status is \texttt{PROPOSAL\_NO\_OBSERVED\_RESULTS}; the observed-results list is empty. Existing clinical findings are cited as evidence anchors only and are not presented as outcomes of the proposed platform.

\balance
\begin{thebibliography}{00}
\bibitem{selkoe2016} D. J. Selkoe and J. Hardy, ``The amyloid hypothesis of Alzheimer's disease at 25 years,'' \emph{EMBO Molecular Medicine}, vol. 8, no. 6, pp. 595--608, 2016, PMID: 27025652.
\bibitem{poewe2017} W. Poewe, K. Seppi, C. M. Tanner, G. M. Halliday, B. Brundin, J. Volkmann, A. E. Schrag, and A. E. Lang, ``Parkinson disease,'' \emph{Nature Reviews Disease Primers}, vol. 3, p. 17013, 2017, PMID: 28332488.
\bibitem{bloem2021} B. R. Bloem, M. S. Okun, and C. Klein, ``Parkinson's disease,'' \emph{The Lancet}, vol. 397, no. 10291, pp. 2284--2303, 2021.
\bibitem{heneka2015} M. T. Heneka et al., ``Neuroinflammation in Alzheimer's disease,'' \emph{The Lancet Neurology}, vol. 14, no. 4, pp. 388--405, 2015, DOI: 10.1016/S1474-4422(15)70016-1.
\bibitem{keren2017} H. Keren-Shaul, A. Spinrad, A. Weiner, O. Matcovitch-Natan, R. Dvir-Szternfeld, T. K. Ulland, E. David, K. Baruch, D. Lara-Astaiso, B. Toth, et al., ``A unique microglia type associated with restricting development of Alzheimer’s disease,'' \emph{Cell}, vol. 169, no. 7, pp. 1276--1290.e17, 2017, PMID: 28602351.
\bibitem{vandyk2023} C. H. van Dyck et al., ``Lecanemab in early Alzheimer's disease,'' \emph{New England Journal of Medicine}, vol. 388, no. 1, pp. 9--21, 2023, PMID: 36449413.
\bibitem{seranova2020} E. Seranova, A. M. Palhegyi, S. Verma, S. Dimova, R. Lasry, M. Naama, C. Sun, T. Barrett, T. R. Rosenstock, D. Kumar, et al., ``Human induced pluripotent stem cell models of neurodegenerative disorders for studying the biomedical implications of autophagy,'' \emph{Journal of Molecular Biology}, vol. 432, no. 8, pp. 2754--2798, 2020, PMID: 32044344.
\bibitem{pushpakom2019} S. Pushpakom, F. Iorio, P. A. Eyers, K. J. Escott, S. Hopper, A. Wells, A. Doig, T. Guilliams, J. Latimer, C. McNamee, et al., ``Drug repurposing: progress, challenges and recommendations,'' \emph{Nature Reviews Drug Discovery}, vol. 18, pp. 41--58, 2019, DOI: 10.1038/nrd.2018.168.
\end{thebibliography}
\end{document}
